%============================================================
% CHƯƠNG 4: THIẾT KẾ, TRIỂN KHAI VÀ ĐÁNH GIÁ HỆ THỐNG
%============================================================
\chapter{THIẾT KẾ, TRIỂN KHAI VÀ ĐÁNH GIÁ HỆ THỐNG}
\label{chap:design}

Chương này trình bày toàn bộ quá trình thiết kế và xây dựng hệ thống
AIoT giám sát và dự báo chất lượng không khí. Nội dung bao gồm thiết
kế kiến trúc tổng thể theo mô hình bốn tầng, thiết kế chi tiết từng
thành phần (phần cứng, firmware, backend, ML pipeline, cơ sở dữ liệu,
dashboard), kết quả xây dựng và kiểm thử, cùng thông tin triển khai
thực tế.

%------------------------------------------------------------
\section{Thiết kế kiến trúc}
\label{sec:design:arch}
%------------------------------------------------------------

\subsection{Lựa chọn kiến trúc phần mềm}
\label{subsec:design:arch:choice}

Hệ thống được thiết kế theo \textbf{kiến trúc bốn tầng} kết hợp mô
hình IoT Gateway, thể hiện trong Hình~\ref{fig:arch_overview}. Kiến
trúc này phân tách rõ ràng trách nhiệm của từng tầng, đảm bảo khả năng
mở rộng độc lập và dễ bảo trì:

\begin{enumerate}[leftmargin=2cm]
  \item \textbf{Tầng thu thập (Edge Layer):} Node cảm biến ESP32/DHT11
    đọc dữ liệu và gửi lên qua MQTT. Tầng này hoạt động hoàn toàn
    tự động, không cần can thiệp người dùng sau khi cấu hình ban đầu.
  \item \textbf{Tầng truyền dẫn (Transport Layer):} Giao thức MQTT
    với Aedes broker tích hợp trong Node.js, đảm bảo truyền dữ liệu
    nhẹ, độ trễ thấp và hỗ trợ many node đồng thời.
  \item \textbf{Tầng xử lý và lưu trữ (Backend Layer):} Node.js/Express
    xử lý nghiệp vụ, PostgreSQL lưu trữ dữ liệu chuỗi thời gian,
    FastAPI/PyTorch chạy LSTM inference và nội suy Kriging/IDW.
  \item \textbf{Tầng hiển thị (Presentation Layer):} React dashboard
    nhận dữ liệu realtime qua Socket.io, hiển thị biểu đồ Chart.js
    và heatmap Leaflet.
\end{enumerate}

\begin{figure}[H]
  \centering
  \includegraphics[width=0.95\textwidth]{arch_overview}
  \caption{Kiến trúc tổng quan của hệ thống AIoT giám sát và dự báo
  chất lượng không khí}
  \label{fig:arch_overview}
\end{figure}

Lý do lựa chọn kiến trúc này so với các phương án khác:

\begin{itemize}[leftmargin=2cm]
  \item \textbf{So với kiến trúc monolithic:} Việc tách FastAPI ML
    service khỏi Node.js backend cho phép mỗi service scale độc lập
    và tránh block event loop khi chạy inference nặng.
  \item \textbf{So với kiến trúc serverless:} Hệ thống cần duy trì
    kết nối MQTT và WebSocket liên tục, không phù hợp với mô hình
    serverless stateless.
  \item \textbf{So với Cloud-only:} Aedes broker tích hợp trong
    Node.js giảm một hop mạng (không cần Mosquitto riêng), giảm độ
    trễ và đơn giản hóa triển khai.
\end{itemize}

\subsection{Thiết kế tổng quan các gói}
\label{subsec:design:arch:packages}

Hình~\ref{fig:package_diagram} mô tả biểu đồ phụ thuộc gói tổng quan
của hệ thống, phân chia thành bốn nhóm chính tương ứng với bốn tầng
kiến trúc.

\begin{figure}[H]
  \centering
  \includegraphics[width=0.95\textwidth]{package_diagram}
  \caption{Biểu đồ phụ thuộc gói tổng quan của hệ thống}
  \label{fig:package_diagram}
\end{figure}

Cấu trúc thư mục dự án phản ánh trực tiếp kiến trúc này:

\begin{itemize}[leftmargin=2cm]
  \item \texttt{sensor/} - Firmware ESP32 (Arduino/C++)
  \item \texttt{server/} - Backend Node.js/Express
    \begin{itemize}
      \item \texttt{controllers/} - Xử lý HTTP request
      \item \texttt{routes/} - Định nghĩa API endpoint
      \item \texttt{services/} - Logic nghiệp vụ
        (alertService, mqttService, forecastService)
      \item \texttt{middleware/} - Xác thực JWT, phân quyền
      \item \texttt{models/} - Kết nối PostgreSQL
    \end{itemize}
  \item \texttt{lstm\_train/} - Python ML pipeline
    (train, inference, spatial\_forecast FastAPI)
  \item \texttt{frontend/src/} - React dashboard
    \begin{itemize}
      \item \texttt{pages/} - Các trang chính
      \item \texttt{components/} - UI components
      \item \texttt{context/} - React Context quản lý
        trạng thái toàn cục
    \end{itemize}
\end{itemize}

%------------------------------------------------------------
\section{Thiết kế chi tiết}
\label{sec:design:detail}
%------------------------------------------------------------

\subsection{Thiết kế phần cứng node cảm biến}
\label{subsec:design:hardware}

\subsubsection*{a. Sơ đồ kết nối phần cứng}

Node cảm biến được xây dựng từ hai thành phần chính: vi điều khiển
ESP32 và cảm biến DHT11. Sơ đồ kết nối được thể hiện trong
Hình~\ref{fig:hardware_schematic}.

\begin{figure}[H]
  \centering
  \includegraphics[width=0.7\textwidth]{hardware_schematic}
  \caption{Sơ đồ kết nối phần cứng node cảm biến}
  \label{fig:hardware_schematic}
\end{figure}

Bảng~\ref{tab:gpio_config} liệt kê cấu hình chân kết nối:

\begin{table}[H]
  \centering
  \caption{Cấu hình chân kết nối ESP32 --- DHT11}
  \label{tab:gpio_config}
  \begin{tabular}{llll}
    \toprule
    \textbf{Chân DHT11} & \textbf{Chân ESP32} & \textbf{Chức năng}
      & \textbf{Ghi chú} \\
    \midrule
    VCC  & 3.3V   & Nguồn cấp      & --- \\
    DATA & GPIO4  & Dữ liệu 1-Wire & Qua điện trở 10k$\Omega$ pull-up \\
    GND  & GND    & Nối đất        & --- \\
    \bottomrule
  \end{tabular}
\end{table}

\subsubsection*{b. Cơ chế lưu cấu hình trên Flash}

Mỗi node cần lưu trữ \texttt{api\_key} và cấu hình nhận được từ
backend qua các lần reset. ESP32 sử dụng thư viện
\texttt{Preferences.h} để ghi vào vùng NVS (Non-Volatile Storage)
của Flash:

\begin{lstlisting}[language=C, caption={Lưu và đọc cấu hình định danh từ Flash ESP32}, label={lst:preferences}]
#include <Preferences.h>
Preferences prefs;

// Đọc device_id và api_key từ NVS Flash khi khởi động
void loadCredentials() {
    prefs.begin("device", true);       // namespace "device", read-only
    deviceId     = prefs.getInt(   "device_id", -1);
    deviceApiKey = prefs.getString("api_key",  "");
    prefs.end();
}

// Ghi device_id và api_key sau khi nhận từ backend qua MQTT
void saveCredentials(int newId, String newKey) {
    prefs.begin("device", false);      // read-write
    prefs.putInt(   "device_id", newId);
    prefs.putString("api_key",   newKey);
    prefs.end();
}
// Lưu ý: logic so sánh ngưỡng cảnh báo (temp_high, hum_low...)
// được thực hiện tập trung ở Backend Node.js (alertService.js),
// không lưu trên firmware để giảm tải cho vi điều khiển.
\end{lstlisting}

\subsection{Thiết kế Firmware}
\label{subsec:design:firmware}

\subsubsection*{a. Kiến trúc firmware phân tầng}

Firmware được tổ chức theo kiến trúc phân tầng như minh họa trong
Hình~\ref{fig:firmware_arch}. Vòng lặp chính (\texttt{loop()}) điều
phối ba tác vụ theo chu kỳ: đọc cảm biến, duy trì kết nối MQTT và
gửi dữ liệu.

\begin{figure}[H]
  \centering
  \includegraphics[width=0.75\textwidth]{firmware_arch}
  \caption{Kiến trúc firmware phân tầng của node ESP32}
  \label{fig:firmware_arch}
\end{figure}

\subsubsection*{b. Luồng xử lý chính của firmware}

\begin{lstlisting}[language=C, caption={Luồng chính firmware ESP32},
label={lst:firmware_main}]
#include <WiFi.h>
#include <PubSubClient.h>
#include <ArduinoJson.h>
#include <DHT.h>
#include <Preferences.h>
#include <WiFiManager.h>

#define DHTPIN 4
#define DHTTYPE DHT11
#define SEND_INTERVAL 5000   // 5 giay

DHT dht(DHTPIN, DHTTYPE);
WiFiClient wifiClient;
PubSubClient mqttClient(wifiClient);

String deviceApiKey = "";
int    deviceId     = -1;
bool   provisioned  = false;

void setup() {
    dht.begin();
    loadCredentials();  // Doc device_id, api_key tu Flash
    WiFiManager wm;
    wm.autoConnect("Sensor-Setup-[MAC]");
    mqttClient.setServer(MQTT_BROKER, 1883);
    mqttClient.setCallback(mqttCallback);
}

void loop() {
    if (!provisioned) { sendProvisionRequest(); return; }
    if (!mqttClient.connected()) reconnectWithAuth();
    mqttClient.loop();

    // Gui du lieu theo chu ky SEND_INTERVAL
    if (millis() - lastSend >= SEND_INTERVAL) {
        float temp = dht.readTemperature();
        float hum  = dht.readHumidity();
        float aqi  = generateMockAqi(); // Giai lap AQI bang Box-Muller

        if (!isnan(temp) && !isnan(hum)) {
            publishData(temp, hum, aqi);
        }
        lastSend = millis();
    }
}

void publishData(float temp, float hum, float aqi) {
    StaticJsonDocument<128> doc;
    doc["temperature"] = temp;
    doc["humidity"]    = hum;
    doc["aqi"]         = aqi;  // bat buoc cho ML pipeline
    char payload[128];
    serializeJson(doc, payload);

    String topic = "device/" + deviceApiKey + "/data";
    mqttClient.publish(topic.c_str(), payload);
}
\end{lstlisting}

\subsubsection*{c. Xử lý message cấu hình từ backend}

Khi nhận cấu hình từ topic \texttt{devices/\{mac\}/config}, callback
MQTT lưu thông tin định danh thiết bị vào Flash và đặt cờ
\texttt{provisioned} để chuyển sang chế độ gửi dữ liệu:

\begin{lstlisting}[language=C, caption={Callback xử lý message cấu
hình MQTT}, label={lst:mqtt_callback}]
void mqttCallback(char* topic, byte* payload,
                  unsigned int length) {
    StaticJsonDocument<256> doc;
    deserializeJson(doc, payload, length);

    // Xu ly lenh reset tu backend (xoa flash, restart)
    if (doc["action"] == "reset") {
        prefs.begin("device", false);
        prefs.clear(); prefs.end();
        ESP.restart(); return;
    }

    // Nhan device_id va api_key tu backend provisioning
    if (doc.containsKey("device_id") &&
        doc.containsKey("api_key")) {
        int    newId  = doc["device_id"].as<int>();
        String newKey = doc["api_key"].as<String>();
        saveCredentials(newId, newKey);
        deviceId     = newId;
        deviceApiKey = newKey;
        provisioned  = true;
    }
    // Logic kiem tra nguong canh bao (temp_high, hum_low...)
    // duoc xu ly tap trung tai alertService.js phia Backend,
    // khong can luu tren firmware.
}
\end{lstlisting}

\subsection{Thiết kế cơ sở dữ liệu}
\label{subsec:design:database}

\subsubsection*{a. Sơ đồ ERD}

Cơ sở dữ liệu PostgreSQL gồm 5 bảng chính, quan hệ được thể hiện
trong Hình~\ref{fig:erd}.

\begin{figure}[H]
  \centering
  \includegraphics[width=\textwidth]{erd}
  \caption{Sơ đồ ERD cơ sở dữ liệu PostgreSQL}
  \label{fig:erd}
\end{figure}

\subsubsection*{b. Định nghĩa schema}

\begin{lstlisting}[language=SQL, caption={Schema cơ sở dữ liệu
PostgreSQL}, label={lst:schema}]
-- Bang nguoi dung
CREATE TABLE users (
    id         SERIAL PRIMARY KEY,
    name       VARCHAR(100),
    email      VARCHAR(100) UNIQUE NOT NULL,
    password   VARCHAR(255) NOT NULL,    -- bcrypt hash
    role       VARCHAR(20) DEFAULT 'user',
    created_at TIMESTAMP DEFAULT NOW()
);
-- Bang phong/khu vuc
CREATE TABLE rooms (
    id            SERIAL PRIMARY KEY,
    name          VARCHAR(100),
    description   TEXT,
    center_lat    REAL,       -- Toa do tam phong
    center_lng    REAL,
    width_m       REAL,       -- Kich thuoc(met)
    length_m      REAL,
    owner_user_id INTEGER REFERENCES users(id),
    created_at    TIMESTAMP DEFAULT NOW()
);
-- Bang thiet bi IoT
CREATE TABLE devices (
    id               SERIAL PRIMARY KEY,
    name             VARCHAR(100),
    description      TEXT,
    owner_user_id    INTEGER REFERENCES users(id),
    api_key          VARCHAR(64) UNIQUE NOT NULL,
    notify_email     BOOLEAN DEFAULT TRUE,
    temp_high        REAL DEFAULT 35.0,
    temp_low         REAL DEFAULT 10.0,
    hum_high         REAL DEFAULT 80.0,
    hum_low          REAL DEFAULT 20.0,
    aqi_high         REAL,
    aqi_low          REAL,
    room_id          INTEGER REFERENCES rooms(id),
    mac_address      VARCHAR(17) UNIQUE,  -- dia chi MAC cua ESP32
    x                DOUBLE PRECISION,   -- toa do X trong phong (m)
    y                DOUBLE PRECISION,   -- toa do Y trong phong (m)
    lat              DOUBLE PRECISION,   -- vi do thuc te (WGS84)
    lng              DOUBLE PRECISION,   -- kinh do thuc te (WGS84)
    firmware_version VARCHAR(50),
    chip_model       VARCHAR(50),
    created_at       TIMESTAMP DEFAULT NOW()
);
-- Bang du lieu cam bien
CREATE TABLE readings (
    id          SERIAL PRIMARY KEY,
    device_id   INTEGER REFERENCES devices(id),
    temperature REAL NOT NULL,
    humidity    REAL NOT NULL,
    aqi         REAL NOT NULL DEFAULT 0,  -- chi so chat luong KK
    timestamp   TIMESTAMP DEFAULT NOW()
);
-- Index tang toc truy van chuoi thoi gian
CREATE INDEX idx_readings_device_time
    ON readings(device_id, timestamp DESC);
-- Bang canh bao
CREATE TABLE alerts (
    id              SERIAL PRIMARY KEY,
    device_id       INTEGER REFERENCES devices(id),
    type            VARCHAR(50),   -- 'temp_high', 'hum_low', 'aqi_high'...
    value           REAL,          -- Gia tri vuot nguong
    message         TEXT,
    acknowledged_by INTEGER REFERENCES users(id),
    acknowledged_at TIMESTAMP,
    timestamp       TIMESTAMP DEFAULT NOW()
);
\end{lstlisting}

\subsubsection*{c. Mối quan hệ giữa các bảng}

\begin{itemize}[leftmargin=2cm]
  \item \texttt{users} $\to$ \texttt{devices}: Một người dùng có thể
    sở hữu nhiều thiết bị. Mỗi thiết bị thuộc một người dùng.
  \item \texttt{users} $\to$ \texttt{rooms}: Một người dùng có thể
    quản lý nhiều phòng.
  \item \texttt{rooms} $\to$ \texttt{devices}: Một phòng có thể chứa
    nhiều node cảm biến. Mỗi node được gán vào một phòng.
  \item \texttt{devices} $\to$ \texttt{readings}: Quan hệ một-nhiều,
    mỗi thiết bị sinh ra nhiều bản ghi theo thời gian.
  \item \texttt{devices} $\to$ \texttt{alerts}: Cảnh báo được tạo tự
    động khi dữ liệu của thiết bị vượt ngưỡng cài đặt.
\end{itemize}

\subsection{Thiết kế Backend API}
\label{subsec:design:backend}

\subsubsection*{a. Cấu trúc API endpoints}

Backend Node.js/Express cung cấp RESTful API với prefix
\texttt{/api}. Tất cả endpoint (trừ auth) yêu cầu JWT trong header
\texttt{Authorization: Bearer \{token\}}.

\begin{table}[H]
  \centering
  \caption{API endpoints quản lý xác thực và người dùng}
  \label{tab:api_auth}
  \begin{tabularx}{\textwidth}{p{1.5cm} p{5cm} X}
    \toprule
    \textbf{Method} & \textbf{Endpoint} & \textbf{Mô tả} \\
    \midrule
    POST & \texttt{/api/auth/register} & Đăng ký tài khoản mới \\
    POST & \texttt{/api/auth/login}    & Đăng nhập, nhận JWT \\
    GET  & \texttt{/api/users/me}      & Lấy thông tin người dùng hiện tại \\
    PUT  & \texttt{/api/users/me}      & Cập nhật thông tin cá nhân \\
    \bottomrule
  \end{tabularx}
\end{table}

\begin{table}[H]
  \centering
  \caption{API endpoints quản lý thiết bị và phòng}
  \label{tab:api_devices}
  \begin{tabularx}{\textwidth}{p{1.5cm} p{4.5cm} X}
    \toprule
    \textbf{Method} & \textbf{Endpoint} & \textbf{Mô tả} \\
    \midrule
    GET    & \texttt{/api/devices}         & Danh sách thiết bị của user \\
    POST   & \texttt{/api/devices}         & Tạo thiết bị mới (Admin) \\
    GET    & \texttt{/api/devices/:id}     & Chi tiết một thiết bị \\
    PUT    & \texttt{/api/devices/:id}     & Cập nhật cấu hình/ngưỡng \\
    DELETE & \texttt{/api/devices/:id}     & Xóa thiết bị \\
    GET    & \texttt{/api/rooms}           & Danh sách phòng \\
    POST   & \texttt{/api/rooms}           & Tạo phòng mới \\
    PUT    & \texttt{/api/rooms/:id}       & Cập nhật thông tin phòng \\
    \bottomrule
  \end{tabularx}
\end{table}

\begin{table}[H]
  \centering
  \caption{API endpoints dữ liệu cảm biến, dự báo và cảnh báo}
  \label{tab:api_data}
  \begin{tabularx}{\textwidth}{p{1.5cm} p{5cm} X}
    \toprule
    \textbf{Method} & \textbf{Endpoint} & \textbf{Mô tả} \\
    \midrule
    GET & \texttt{/api/data}              & Lịch sử dữ liệu theo \texttt{device\_id}, \texttt{start}, \texttt{end} \\
    GET & \texttt{/api/forecast}          & Gọi FastAPI LSTM, trả về 15 giá trị dự báo \\
    GET & \texttt{/api/spatial-forecast}  & Gọi FastAPI Kriging, trả về lưới 40$\times$50 \\
    GET & \texttt{/api/alerts}            & Danh sách cảnh báo (lọc theo device, thời gian) \\
    PUT & \texttt{/api/alerts/:id/ack}    & Xác nhận (acknowledge) một cảnh báo \\
    \bottomrule
  \end{tabularx}
\end{table}

\subsubsection*{b. Cơ chế xác thực JWT}

\begin{lstlisting}[language=JavaScript, caption={Middleware xác thực
JWT}, label={lst:jwt_middleware}]
// middleware/auth.js
const jwt = require('jsonwebtoken');
module.exports = (req, res, next) => {
    const authHeader = req.headers['authorization'];
    if (!authHeader) return res.status(401)
        .json({ error: 'Missing token' });
    const token = authHeader.split(' ')[1]; // Bearer <token>
    try {
        const decoded = jwt.verify(token, process.env.JWT_SECRET);
        req.user = decoded;   // { id, email, role }
        next();
    } catch (err) {
        return res.status(403).json({ error: 'Invalid token' });
    }
};
\end{lstlisting}

\subsubsection*{c. Dịch vụ cảnh báo (alertService)}

\begin{lstlisting}[language=JavaScript, caption={Logic kiem tra nguong canh bao}, label={lst:alert_service}]
// services/alertService.js
async function check(device, reading) {
    const { temperature, humidity } = reading;
    const alerts = [];

    if (temperature > device.temp_high)
        alerts.push({ type: 'temp_high', value: temperature,
            message: `Nhiet do ${temperature}C vuot nguong ${device.temp_high}C` });
            
    if (temperature < device.temp_low)
        alerts.push({ type: 'temp_low', value: temperature,
            message: `Nhiet do ${temperature}C duoi nguong ${device.temp_low}C` });
            
    if (humidity > device.hum_high)
        alerts.push({ type: 'hum_high', value: humidity,
            message: `Do am ${humidity}% vuot nguong ${device.hum_high}%` });
            
    if (humidity < device.hum_low)
        alerts.push({ type: 'hum_low', value: humidity,
            message: `Do am ${humidity}% duoi nguong ${device.hum_low}%` });
            
    for (const alert of alerts) {
        // Luu vao DB
        await db.query(
            `INSERT INTO alerts
             (device_id, type, value, message, timestamp)
             VALUES ($1,$2,$3,$4,NOW())`,
            [device.id, alert.type,
             alert.value, alert.message]);
             
        // Gui email neu thiet bi bat notify
        if (device.notify_email)
            await emailService.sendAlert(device, alert);
    }
}
\end{lstlisting}

\subsection{Thiết kế và huấn luyện mô hình Machine Learning}
\label{subsec:design:ml}

\subsubsection*{a. Chuẩn bị dữ liệu}

Dữ liệu huấn luyện được thu thập từ hệ thống trong quá trình vận hành
thử nghiệm, lưu trữ trong bảng \texttt{readings}. Quy trình chuẩn bị
dữ liệu gồm các bước:

\begin{enumerate}[leftmargin=2cm]
  \item \textbf{Truy vấn:} Lấy toàn bộ dữ liệu từ PostgreSQL theo
    thứ tự thời gian tăng dần.
  \item \textbf{Kiểm tra missing values:} Bỏ qua các bản ghi có giá
    trị \texttt{NULL} hoặc nằm ngoài dải hợp lệ (nhiệt độ: $-10$ đến
    $60^\circ$C, độ ẩm: $0$ đến $100$\%).
  \item \textbf{Chuẩn hóa:} Áp dụng \texttt{MinMaxScaler}
    (scikit-learn) để đưa tất cả giá trị về khoảng $[0, 1]$. Lưu
    scaler vào \texttt{scaler.pkl} để tái dùng khi inference.
  \item \textbf{Tạo chuỗi (sliding window):} Với mỗi vị trí $t$,
    tạo cặp $(X_t, y_t)$ trong đó $X_t$ là ma trận $30 \times 3$
    (30 bước quá khứ, 3 features: nhiệt độ, độ ẩm, AQI) và $y_t$ là ma trận $15 \times 3$
    (15 bước tương lai cần dự báo).
  \item \textbf{Chia tập:} 80\% dữ liệu cho huấn luyện, 20\% cho
    kiểm thử.
\end{enumerate}

\subsubsection*{b. Kiến trúc mô hình LSTM}

Bảng~\ref{tab:lstm_arch} mô tả kiến trúc mô hình LSTM được sử dụng
trong đồ án.

\begin{table}[H]
  \centering
  \caption{Kiến trúc mô hình LSTM}
  \label{tab:lstm_arch}
  \begin{tabular}{lllr}
    \toprule
    \textbf{Layer} & \textbf{Loại} & \textbf{Tham số} &
    \textbf{Output shape} \\
    \midrule
    Input   & ---          & ---
      & $(batch, 30, 3)$ \\
    LSTM 1  & LSTM         & hidden=64, dropout=0.3
      & $(batch, 30, 64)$ \\
    LSTM 2  & LSTM         & hidden=64, dropout=0.3
      & $(batch, 64)$ \\
    Linear  & Fully connected & in=64, out=$15 \times 3$
      & $(batch, 45)$ \\
    Reshape & ---          & ---
      & $(batch, 15, 3)$ \\
    \bottomrule
    \multicolumn{3}{l}{\textit{Tổng số tham số:}}
      & \textit{$\approx$68.000} \\
  \end{tabular}
\end{table}

\subsubsection*{c. Huấn luyện và đánh giá}

Mô hình được huấn luyện với các siêu tham số:

\begin{table}[H]
  \centering
  \caption{Siêu tham số huấn luyện mô hình LSTM}
  \label{tab:lstm_hyperparams}
  \begin{tabular}{ll}
    \toprule
    \textbf{Siêu tham số} & \textbf{Giá trị} \\
    \midrule
    Optimizer      & Adam ($lr = 10^{-3}$) \\
    Loss function  & MSELoss \\
    Batch size     & 32 \\
    Số epoch tối đa& 100 \\
    Early stopping & patience = 10 epoch \\
    Dropout        & 0.3 (áp dụng giữa các LSTM layer) \\
    \bottomrule
  \end{tabular}
\end{table}

Kết quả huấn luyện thể hiện qua loss curve lưu tại
\texttt{lstm\_train/loss\_curve.png} (ngắn hạn) và
\texttt{loss\_curve\_long.png} (dài hạn). Đồ thị cho thấy training
loss và validation loss hội tụ ổn định, không có hiện tượng overfitting
rõ rệt nhờ cơ chế dropout và early stopping.

\subsection{Thiết kế Dashboard (Frontend)}
\label{subsec:design:frontend}

\subsubsection*{a. Sơ đồ điều hướng trang}

Hình~\ref{fig:nav_diagram} mô tả sơ đồ điều hướng giữa các trang
trong React dashboard. Trang Login là điểm vào bắt buộc; sau khi xác
thực thành công, người dùng được chuyển đến Dashboard chính.

\begin{figure}[H]
  \centering
  \includegraphics[width=0.85\textwidth]{nav_diagram}
  \caption{Sơ đồ điều hướng các trang trong React dashboard}
  \label{fig:nav_diagram}
\end{figure}

Các trang chính trong ứng dụng:

\begin{itemize}[leftmargin=2cm]
  \item \textbf{Login / Register:} Xác thực người dùng, lưu JWT vào
    localStorage.
  \item \textbf{Dashboard (trang chủ):} Hiển thị tổng quan tất cả
    node, biểu đồ realtime nhiệt độ, độ ẩm và AQI cập nhật qua Socket.io.
  \item \textbf{Trang thiết bị:} Danh sách node, cấu hình ngưỡng
    cảnh báo, xem biểu đồ lịch sử và dự báo LSTM cho từng node.
  \item \textbf{Trang dự báo (Forecast):} Gồm hai chế độ xem được
    chuyển đổi bằng tab: (1) \textbf{Single Device} --- biểu đồ dự
    báo chuỗi thời gian LSTM cho từng node theo các mốc 6/10/15 phút
    (ngắn hạn) và 1/3/6/12/24 giờ (dài hạn); (2) \textbf{Spatial} ---
    bản đồ heatmap Kriging/IDW phân bố môi trường không gian theo
    phòng, render trên Leaflet.
  \item \textbf{Trang cảnh báo:} Danh sách alert, bộ lọc theo thời
    gian, chức năng acknowledge.
  \item \textbf{Trang quản lý phòng:} Tạo và cấu hình phòng, gán
    node vào phòng, cài đặt tọa độ.
\end{itemize}

\subsubsection*{b. Thiết kế giao diện người dùng}

Hình~\ref{fig:ui_dashboard} và Hình~\ref{fig:ui_heatmap} minh họa
thiết kế giao diện của hai màn hình chính.

\begin{figure}[H]
  \centering
  \includegraphics[width=\textwidth]{ui_dashboard}
  \caption{Giao diện Dashboard chính hiển thị dữ liệu realtime}
  \label{fig:ui_dashboard}
\end{figure}

\begin{figure}[H]
  \centering
  \includegraphics[width=\textwidth]{system_spatial_forecast}
  \caption{Giao diện chức năng dự báo không gian (Spatial Forecast)}
  \label{fig:ui_heatmap}
\end{figure}

%------------------------------------------------------------
\section{Xây dựng hệ thống}
\label{sec:design:impl}
%------------------------------------------------------------

\subsection{Thư viện và công cụ sử dụng}
\label{subsec:design:tools}

\begin{table}[H]
  \centering
  \caption{Danh sách thư viện và công cụ sử dụng}
  \label{tab:tools}
  \begin{tabularx}{\textwidth}{p{2.5cm} p{3.5cm} p{2cm} X}
    \toprule
    \textbf{Mục đích} & \textbf{Thư viện/Công cụ}
      & \textbf{Phiên bản} & \textbf{Ghi chú} \\
    \midrule
    Firmware IDE  & Arduino IDE     & 2.x     & Biên dịch C++ ESP32 \\
    MQTT Client   & PubSubClient    & 2.8     & MQTT trên ESP32 \\
    JSON (FW)     & ArduinoJson     & 6.x     & Đóng gói payload \\
    Cảm biến      & DHT sensor lib  & 1.4     & Đọc DHT11 \\
    Backend FW    & Node.js         & 18 LTS  & Runtime JavaScript \\
    Web framework & Express         & 4.x     & REST API \\
    MQTT Broker   & Aedes           & 0.50    & Broker tích hợp \\
    Realtime      & Socket.io       & 4.x     & WebSocket server \\
    Database      & PostgreSQL      & 15      & RDBMS chính \\
    ORM/Query     & node-postgres   & 8.x     & PostgreSQL client \\
    Auth          & jsonwebtoken    & 9.x     & JWT \\
    Password      & bcrypt          & 5.x     & Hash mật khẩu \\
    Email         & nodemailer      & 6.x     & Gửi cảnh báo \\
    ML Framework  & PyTorch         & 2.x     & Huấn luyện LSTM \\
    ML API        & FastAPI         & 0.110   & Serve ML inference \\
    Nội suy       & pykrige         & 1.7     & Kriging \\
    Nội suy       & scipy           & 1.12    & IDW \\
    Scaler        & scikit-learn    & 1.4     & MinMaxScaler \\
    Frontend      & React           & 19      & UI framework \\
    Build tool    & Vite            & 5.x     & Bundler \\
    Biểu đồ       & Chart.js        & 4.x     & Vẽ biểu đồ \\
    Bản đồ        & Leaflet         & 1.9     & Render heatmap \\
    Quản lý mã    & Git / GitHub    & ---     & Version control \\
    \bottomrule
  \end{tabularx}
\end{table}

\subsection{Kết quả đạt được}
\label{subsec:design:result}

Sau quá trình thiết kế và triển khai, đồ án đã xây dựng thành công hệ
thống hoàn chỉnh với đầy đủ các thành phần hoạt động ổn định. Kết quả
cụ thể từng thành phần:

\begin{itemize}[leftmargin=2cm]
  \item \textbf{Node cảm biến ESP32:} Hoạt động ổn định, đọc và gửi
    dữ liệu định kỳ 5 giây qua MQTT. Cơ chế tự động reconnect đảm bảo
    node phục hồi kết nối sau khi mất mạng mà không cần can thiệp thủ
    công.
  \item \textbf{Backend Node.js:} Tiếp nhận và xử lý dữ liệu từ nhiều
    node đồng thời. MQTT broker Aedes và Socket.io hoạt động ổn định
    trong cùng tiến trình. API trả về kết quả trong thời gian trung
    bình dưới 100ms cho các truy vấn dữ liệu thông thường.
  \item \textbf{FastAPI ML service:} Inference LSTM hoàn thành trong
    dưới 200ms cho một node (window 30, horizon 15). Kriging tính toán
    lưới 40$\times$50 trong khoảng 1--2 giây tùy số lượng node.
  \item \textbf{Dashboard React:} Biểu đồ Chart.js cập nhật realtime
    khi nhận event Socket.io, không cần tải lại trang. Heatmap Leaflet
    render đúng và trực quan trên các trình duyệt hiện đại.
  \item \textbf{Mô hình LSTM:} Huấn luyện hội tụ tốt trên tập dữ liệu
    thu thập được. Loss curve thể hiện xu hướng giảm đều và ổn định
    giữa training loss và validation loss.
\end{itemize}

Một số ảnh tiêu biểu được chọn để chèn trực tiếp trong nội dung báo cáo, minh họa toàn bộ các giao diện chức năng của hệ thống[cite: 122, 136].

% CỤM 1: Xác thực & Quản lý thiết bị (STT 1 - 4)
\begin{figure}[H]
  \centering
  \begin{subfigure}{0.48\textwidth}
    \centering
    \includegraphics[width=\textwidth]{system_login.jpg}
    \caption{Màn hình đăng nhập}
  \end{subfigure}
  \hfill
  \begin{subfigure}{0.48\textwidth}
    \centering
    \includegraphics[width=\textwidth]{system_register.jpg}
    \caption{Màn hình đăng ký}
  \end{subfigure}

  \vspace{0.4cm}
  \begin{subfigure}{0.48\textwidth}
    \centering
    \includegraphics[width=\textwidth]{system_devices.jpg}
    \caption{Danh sách thiết bị IoT}
  \end{subfigure}
  \hfill
  \begin{subfigure}{0.48\textwidth}
    \centering
    \includegraphics[width=\textwidth]{system_devices_setting.jpg}
    \caption{Cấu hình ngưỡng cảnh báo}
  \end{subfigure}
  \caption{Giao diện chức năng hệ thống --- Đăng nhập, đăng ký và quản lý thiết bị}
  \label{fig:system_demo_part1}
\end{figure}

% CỤM 2: Dashboard & Chi tiết thiết bị (STT 5 - 8)
\begin{figure}[H]
  \centering
  \begin{subfigure}{0.48\textwidth}
    \centering
    \includegraphics[width=\textwidth]{system_dashboard.jpg}
    \caption{Dashboard tổng quan}
  \end{subfigure}
  \hfill
  \begin{subfigure}{0.48\textwidth}
    \centering
    \includegraphics[width=\textwidth]{system_devices_detail_chart.jpg}
    \caption{Biểu đồ chi tiết thiết bị (Live Chart)}
  \end{subfigure}

  \vspace{0.4cm}
  \begin{subfigure}{0.48\textwidth}
    \centering
    \includegraphics[width=\textwidth]{system_devices_detail_history.jpg}
    \caption{Lịch sử dữ liệu cảm biến}
  \end{subfigure}
  \hfill
  \begin{subfigure}{0.48\textwidth}
    \centering
    \includegraphics[width=\textwidth]{system_devices_detail_alert.jpg}
    \caption{Lịch sử cảnh báo của thiết bị}
  \end{subfigure}
  \caption{Giao diện chức năng hệ thống --- Dashboard và các tab chi tiết thiết bị}
  \label{fig:system_demo_part2}
\end{figure}

% CỤM 3: Dự báo & Quản lý không gian, Cảnh báo hệ thống (STT 9 - 13)
\begin{figure}[H]
  \centering
  \begin{subfigure}{0.48\textwidth}
    \centering
    \includegraphics[width=\textwidth]{system_forecast.jpg}
    \caption{Dự báo chuỗi thời gian LSTM}
  \end{subfigure}
  \hfill
  \begin{subfigure}{0.48\textwidth}
    \centering
    \includegraphics[width=\textwidth]{system_spatial_forecast.jpg}
    \caption{Heatmap không gian}
  \end{subfigure}

  \vspace{0.4cm}
  \begin{subfigure}{0.48\textwidth}
    \centering
    \includegraphics[width=\textwidth]{system_rooms.jpg}
    \caption{Quản lý danh sách phòng}
  \end{subfigure}
  \hfill
  \begin{subfigure}{0.48\textwidth}
    \centering
    \includegraphics[width=\textwidth]{system_rooms_config.jpg}
    \caption{Cấu hình kích thước và vị trí node}
  \end{subfigure}

  \vspace{0.4cm}
  \begin{subfigure}{0.48\textwidth}
    \centering
    \includegraphics[width=\textwidth]{system_alerts.jpg}
    \caption{Danh sách cảnh báo hệ thống}
  \end{subfigure}
  \caption{Giao diện dự báo thời gian/không gian, quản lý phòng và cảnh báo hệ thống}
  \label{fig:system_demo_part3}
\end{figure}

\subsection{Minh họa các chức năng chính}
\label{subsec:design:demo}

\textbf{Chức năng 1 --- Đăng ký node IoT tự động:} Khi node ESP32
khởi động lần đầu, hệ thống tự động cấp phát \texttt{api\_key} qua
MQTT mà không cần thao tác thủ công từ phía người dùng. Node sẵn sàng
gửi dữ liệu trong vòng vài giây sau khi kết nối Wi-Fi thành công.

\textbf{Chức năng 2 --- Giám sát realtime:} Biểu đồ nhiệt độ và độ
ẩm trên dashboard cập nhật tức thì mỗi khi node gửi dữ liệu mới qua
Socket.io, không cần tải lại trang. Người dùng có thể theo dõi nhiều
node cùng lúc trên cùng giao diện.

\textbf{Chức năng 3 --- Dự báo chuỗi thời gian:} Người dùng chọn node
và khoảng thời gian, hệ thống trả về biểu đồ dự báo 15 bước tiếp theo
(màu khác để phân biệt với dữ liệu lịch sử). Thời gian phản hồi từ
khi nhấn nút đến khi hiện kết quả dưới 3 giây.

\textbf{Chức năng 4 --- Heatmap phân bố không gian:} Từ dữ liệu của
các node trong phòng, hệ thống render heatmap nhiệt độ/độ ẩm phủ lên
bản đồ sàn phòng. Gradient màu từ xanh lam (mát/khô) đến đỏ
(nóng/ẩm) giúp người dùng nhận biết vùng bất thường trực quan.

\textbf{Chức năng 5 --- Cảnh báo tự động:} Khi nhiệt độ hoặc độ ẩm
vượt ngưỡng cài đặt, hệ thống tự động ghi alert, hiển thị thông báo
trên dashboard và gửi email đến người dùng trong vòng 30 giây.

%------------------------------------------------------------
\section{Kiểm thử}
\label{sec:design:test}
%------------------------------------------------------------

Kiểm thử được thực hiện theo phương pháp hộp đen (black-box), tập
trung vào đầu vào và đầu ra quan sát được mà không xét cấu trúc nội
bộ mã nguồn.

\begin{table}[H]
  \centering
  \caption{Kiểm thử chức năng đăng ký và kết nối node IoT}
  \label{tab:test_register}
  \begin{tabularx}{\textwidth}{p{4cm} X p{1.5cm}}
    \toprule
    \textbf{Trường hợp kiểm thử} & \textbf{Kết quả mong đợi}
      & \textbf{Kết quả} \\
    \midrule
    Node mới kết nối lần đầu
      & Backend tạo bản ghi device, sinh \texttt{api\_key}, gửi config
        về node qua MQTT
      & Đạt \\
    Node đã tồn tại kết nối lại
      & Backend trả về cấu hình hiện có, không tạo bản ghi mới
      & Đạt \\
    Wi-Fi mất kết nối giữa chừng
      & Node tự động reconnect sau 5 giây
      & Đạt \\
    \texttt{api\_key} không hợp lệ
      & Backend bỏ qua message, ghi log cảnh báo
      & Đạt \\
    \bottomrule
  \end{tabularx}
\end{table}

\begin{table}[H]
  \centering
  \caption{Kiểm thử chức năng thu thập và hiển thị dữ liệu realtime}
  \label{tab:test_realtime}
  \begin{tabularx}{\textwidth}{p{4cm} X p{1.5cm}}
    \toprule
    \textbf{Trường hợp kiểm thử} & \textbf{Kết quả mong đợi}
      & \textbf{Kết quả} \\
    \midrule
    Node gửi dữ liệu bình thường
      & Dashboard cập nhật biểu đồ trong dưới 2 giây
      & Đạt \\
    DHT11 trả về NaN (lỗi đọc)
      & Firmware bỏ qua chu kỳ đó, thử lại sau 5 giây
      & Đạt \\
    Nhiều node gửi đồng thời
      & Tất cả dữ liệu được lưu đúng, không bị mất
      & Đạt \\
    Client mở nhiều tab cùng lúc
      & Mỗi tab nhận đúng dữ liệu của node đang subscribe
      & Đạt \\
    \bottomrule
  \end{tabularx}
\end{table}

\begin{table}[H]
  \centering
  \caption{Kiểm thử chức năng dự báo và heatmap}
  \label{tab:test_forecast}
  \begin{tabularx}{\textwidth}{p{4cm} X p{1.5cm}}
    \toprule
    \textbf{Trường hợp kiểm thử} & \textbf{Kết quả mong đợi}
      & \textbf{Kết quả} \\
    \midrule
    Yêu cầu dự báo khi đủ 30 bản ghi
      & FastAPI trả về mảng 15 giá trị, dashboard hiển thị đúng
      & Đạt \\
    Yêu cầu dự báo khi chưa đủ 30 bản ghi
      & API trả về lỗi 400 kèm thông báo rõ ràng
      & Đạt \\
    FastAPI service bị dừng
      & Backend Node.js trả về lỗi 503, không crash
      & Đạt \\
    Heatmap với 1 node duy nhất
      & Hệ thống sinh node ảo bù biên, render được heatmap
      & Đạt \\
    \bottomrule
  \end{tabularx}
\end{table}

\begin{table}[H]
  \centering
  \caption{Kiểm thử chức năng cảnh báo}
  \label{tab:test_alert}
  \begin{tabularx}{\textwidth}{p{4cm} X p{1.5cm}}
    \toprule
    \textbf{Trường hợp kiểm thử} & \textbf{Kết quả mong đợi}
      & \textbf{Kết quả} \\
    \midrule
    Nhiệt độ vượt ngưỡng cao
      & Alert được tạo trong DB, email gửi đến người dùng
      & Đạt \\
    Độ ẩm dưới ngưỡng thấp
      & Alert được tạo, hiển thị trên dashboard
      & Đạt \\
    Xác nhận (acknowledge) alert
      & Alert được đánh dấu đã xử lý, không gửi email lần nữa
      & Đạt \\
    Tắt tính năng notify\_email
      & Alert vẫn ghi DB nhưng không gửi email
      & Đạt \\
    \bottomrule
  \end{tabularx}
\end{table}

%------------------------------------------------------------
\section{Triển khai}
\label{sec:design:deploy}
%------------------------------------------------------------

\subsection{Môi trường triển khai}
\label{subsec:design:deploy:env}

Hệ thống được triển khai và kiểm thử trong môi trường phòng lab với
cấu hình cụ thể như sau:

\begin{table}[H]
  \centering
  \caption{Cấu hình môi trường triển khai}
  \label{tab:deploy_env}
  \begin{tabular}{ll}
    \toprule
    \textbf{Thành phần} & \textbf{Cấu hình} \\
    \midrule
    Máy chủ Backend     & Máy tính cá nhân / Laptop \\
    Hệ điều hành        & Ubuntu 22.04 LTS / Windows 11 (WSL2) \\
    Node.js             & v18 LTS \\
    Python              & 3.10+ \\
    PostgreSQL          & v15, chạy local \\
    Node IoT            & ESP32 DevKit v1 + DHT11 \\
    Mạng                & Wi-Fi LAN cùng subnet \\
    \bottomrule
  \end{tabular}
\end{table}

\subsection{Hướng dẫn khởi động hệ thống}
\label{subsec:design:deploy:guide}

Quy trình khởi động hệ thống theo thứ tự:

\begin{enumerate}[leftmargin=2cm]
  \item \textbf{Khởi động PostgreSQL} và chạy script tạo bảng
    (\texttt{schema.sql}).
  \item \textbf{Khởi động Backend Node.js:}
\begin{lstlisting}[language=bash, label={lst:start_backend}]
cd server && npm install
cp .env.example .env   # Cau hinh DB, JWT_SECRET, email
node index.js
\end{lstlisting}

  \item \textbf{Khởi động hai FastAPI ML service (chạy đồng thời):}
\begin{lstlisting}[language=bash, label={lst:start_fastapi}]
cd forecast
# Service 1: LSTM time-series forecast (cong 8000)
cd time_forecast
uvicorn main:app --port 8000 --reload &

# Service 2: LSTM + Kriging spatial forecast (cong 8001)
cd ..
uvicorn spatial_forecast:app --port 8001 --reload
\end{lstlisting}

  \item \textbf{Khởi động Frontend:}
\begin{lstlisting}[language=bash, label={lst:start_frontend}]
cd frontend && npm install
npm run dev   # Development
npm run build && npm run preview   # Production
\end{lstlisting}

  \item \textbf{Nạp firmware lên ESP32:} Mở
    \texttt{sensor/sensor.ino} trong Arduino IDE, cấu hình
    \texttt{WIFI\_SSID}, \texttt{WIFI\_PASSWORD} và
    \texttt{MQTT\_BROKER} (địa chỉ IP máy chủ), nạp firmware.
  \item \textbf{Kiểm tra:} Mở trình duyệt tại
    \texttt{http://localhost:5173}, đăng ký tài khoản, xem node
    tự động xuất hiện trong danh sách thiết bị.
\end{enumerate}

\subsection{Script mô phỏng node IoT}
\label{subsec:design:deploy:simulate}

Để kiểm thử hệ thống mà không cần phần cứng thực, đồ án cung cấp
script mô phỏng \texttt{scripts/simulate\_esp32.js}:

\begin{lstlisting}[language=JavaScript, caption={Script mô phỏng node
ESP32 bằng Node.js}, label={lst:simulate}]
// scripts/simulate_esp32.js
const mqtt = require('mqtt');
const client = mqtt.connect('mqtt://localhost:1883');

const API_KEY = 'your-api-key-here';
const TOPIC   = `device/${API_KEY}/data`;
setInterval(() => {
    // Sinh dữ liệu ngẫu nhiên trong dải thực tế
    const payload = JSON.stringify({
        temperature: (20 + Math.random() * 15).toFixed(1),
        humidity:    (40 + Math.random() * 40).toFixed(1),
        aqi:         (30 + Math.random() * 70).toFixed(1)  // bat buoc cho ML
    });
    client.publish(TOPIC, payload);
    console.log(`Published: ${payload}`);
}, 5000);
\end{lstlisting}

Script này cho phép kiểm thử toàn bộ pipeline từ MQTT đến dashboard
mà không cần thiết bị ESP32 vật lý.